% !TeX root = ../抽象代数笔记.tex
% 由主文档通过 \input 载入；请编译主文档。
\setcounter{chapter}{3} % 保留原书章号。
\chapter{陪集、指数与拉格朗日定理}\label{ch:cosets}

\begin{description}
  \item[学习目标] 用陪集划分理解整除结论，区分集合的分块与商群的运算。
  \item[先修知识] 第\ref{ch:groups}、\ref{ch:cyclic-permutation}章。
\end{description}

\section{左陪集、右陪集与等价关系}
对 $H\leq G$ 及 $g\in G$ 定义左陪集 $gH$、右陪集 $Hg$。证明左陪集相交即相等，并且每个左陪集与 $H$ 等势。定义指数 $[G:H]$ 为左陪集的个数，说明左、右陪集数量相同，但相应陪集不必相等。

\section{拉格朗日定理及其边界}
对有限群 $G$ 与子群 $H$，证明拉格朗日定理（Lagrange's theorem）$|G|=[G:H]|H|$。推出元素阶整除群阶、素数阶群必为循环群。以 $A_4$ 没有 $6$ 阶子群作为反向命题失败的例子，其完整证明放在正规子群工具建立后完成。

\section{指数的乘法公式}
在 $K\leq H\leq G$ 且相关指数有限时，证明 $[G:K]=[G:H][H:K]$。通过选择陪集代表元给出可核查的计数论证。说明陪集代表元通常不唯一，计数结果却不依赖代表元的具体选择。

\section{数论应用与例题练习}
在模 $n$ 可逆剩余类构成的乘法群中推导欧拉定理，并以素数模数得到费马小定理。计算 $S_3$ 中不同子群的左右陪集，比较陪集数量与子群个数。自检要求解释为什么仅有 $|H|\mid |G|$ 不能保证某个指定子集是子群，也不能保证每个因子都对应子群。

